% GENERATED FILE. Edit canonical lessons, then run npm run generate:books.
% canonical-source-hash: fnv1a64:a1933a640de471af
% canonical-lessons: FA-C347-afsane, FA-C347-osture, FA-C347-qahreman, FA-C347-ezhdeha, FA-C347-div, FA-R347-review-340-347

\chapter{Legends and heroes}
\label{ch:fa-legends-and-heroes}
\hlchaptermodality{347}

\begin{chapteropening}
\textbf{By the end of this chapter:} \emph{I can talk about legends, myths, heroes, dragons and demons.}

Complete the last lesson of chapter 347: I can talk about legends, myths, heroes, dragons and demons.
\end{chapteropening}

\section[\texorpdfstring{\fa{افسانه} (afsâne)}{افسانه (afsâne)}]{\fa{افسانه} (afsâne) — a legend, a fable}

\label{lesson:FA-C347-afsane}



\begin{quote}
Before the new one: say the Persian for to relate, then the Persian for to remind.
\end{quote}

\subsection*{You'll want to know: \fa{افسانه}}
\textbf{\fa{افسانه}} (\emph{afsâne}) — \textquotedblleft{}a legend, a fable\textquotedblright{}.

\begin{quote}
\fa{این} \fa{افسانه} \fa{درباره} \fa{یک} \fa{کوه} \fa{بلند} \fa{است}.
\end{quote}

\begin{quote}
\emph{In afsâne darbâre-ye yek kuh-e boland ast.}
\end{quote}

\begin{quote}
This legend is about a high mountain.
\end{quote}

Also \textquotedblleft{}a made-up story\textquotedblright{}: \emph{afsâne ast} \textquotedblleft{}it is a fairy tale\textquotedblright{}. \emph{Afsâne-i} \textquotedblleft{}legendary\textquotedblright{}.

\begin{cousinweb}
\emph{afsâne} is inherited from Middle Persian \emph{afsānag}, \textquotedblleft{}fable\textquotedblright{} (EtymDB).
\end{cousinweb}

\begin{grammarlens}[title={What you've built}]
One more word for telling, explaining and describing at B1.
\end{grammarlens}

\subsection*{Your turn}
\begin{itemize}
\raggedright
  \item \emph{Say it:} \emph{afsâne}
  \item \emph{Say it:} \emph{afsâne}, once more
  \item \emph{Say it:} a sentence of your own with \emph{afsâne}
  \item \emph{Recall:} say the Persian for to relate, then the Persian for to remind, then say \emph{afsâne} again
\end{itemize}

\begin{tcolorbox}[breakable,colback=teal!4,colframe=teal!35!black,title={Before you move on}]
What does \fa{افسانه} mean? (\textquotedblleft{}A legend, a fable\textquotedblright{}.) Say it once more.
\end{tcolorbox}
\section[\texorpdfstring{\fa{اسطوره} (osture)}{اسطوره (osture)}]{\fa{اسطوره} (osture) — a myth}

\label{lesson:FA-C347-osture}



\begin{quote}
Before the new one: say the Persian for to remind, then the Persian for a legend.
\end{quote}

\subsection*{You'll want to know: \fa{اسطوره}}
\textbf{\fa{اسطوره}} (\emph{osture}) — \textquotedblleft{}a myth\textquotedblright{}.

\begin{quote}
\fa{در} \fa{این} \fa{اسطوره} \fa{آتش} \fa{مقدس} \fa{است}.
\end{quote}

\begin{quote}
\emph{Dar in osture âtash moqaddas ast.}
\end{quote}

\begin{quote}
In this myth fire is sacred.
\end{quote}

Broken plural \emph{asâtir}. The word reached Arabic from Greek.

\begin{cousinweb}
\emph{osture} is borrowed from Arabic \emph{usṭūra}, \textquotedblleft{}fable, myth\textquotedblright{}, and EtymDB takes the Arabic word from Greek \emph{historía}, \textquotedblleft{}history, inquiry\textquotedblright{}.
\end{cousinweb}

\begin{grammarlens}[title={What you've built}]
One more word for telling, explaining and describing at B1.
\end{grammarlens}

\subsection*{Your turn}
\begin{itemize}
\raggedright
  \item \emph{Say it:} \emph{osture}
  \item \emph{Say it:} \emph{osture}, once more
  \item \emph{Say it:} a sentence of your own with \emph{osture}
  \item \emph{Recall:} say the Persian for to remind, then the Persian for a legend, then say \emph{osture} again
\end{itemize}

\begin{tcolorbox}[breakable,colback=teal!4,colframe=teal!35!black,title={Before you move on}]
What does \fa{اسطوره} mean? (\textquotedblleft{}A myth\textquotedblright{}.) Say it once more.
\end{tcolorbox}
\section[\texorpdfstring{\fa{قهرمان} (qahremân)}{قهرمان (qahremân)}]{\fa{قهرمان} (qahremân) — a hero; a champion}

\label{lesson:FA-C347-qahreman}



\begin{quote}
Before the new one: say the Persian for a legend, then the Persian for a myth.
\end{quote}

\subsection*{You'll want to know: \fa{قهرمان}}
\textbf{\fa{قهرمان}} (\emph{qahremân}) — \textquotedblleft{}a hero; a champion\textquotedblright{}.

\begin{quote}
\fa{قهرمان} \fa{داستان} \fa{یک} \fa{پسر} \fa{فقیر} \fa{است}.
\end{quote}

\begin{quote}
\emph{Qahremân-e dâstân yek pesar-e faqir ast.}
\end{quote}

\begin{quote}
The hero of the story is a poor boy.
\end{quote}

Also a sports champion; \emph{qahremâni} \textquotedblleft{}a championship\textquotedblright{}.

\begin{cousinweb}
EtymDB glosses \emph{qahremân} \textquotedblleft{}champion, hero\textquotedblright{} and gives no parent; its origin is not settled.
\end{cousinweb}

\begin{grammarlens}[title={What you've built}]
One more word for telling, explaining and describing at B1.
\end{grammarlens}

\subsection*{Your turn}
\begin{itemize}
\raggedright
  \item \emph{Say it:} \emph{qahremân}
  \item \emph{Say it:} \emph{qahremân}, once more
  \item \emph{Say it:} a sentence of your own with \emph{qahremân}
  \item \emph{Recall:} say the Persian for a legend, then the Persian for a myth, then say \emph{qahremân} again
\end{itemize}

\begin{tcolorbox}[breakable,colback=teal!4,colframe=teal!35!black,title={Before you move on}]
What does \fa{قهرمان} mean? (\textquotedblleft{}A hero; a champion\textquotedblright{}.) Say it once more.
\end{tcolorbox}
\section[\texorpdfstring{\fa{اژدها} (ezhdehâ)}{اژدها (ezhdehâ)}]{\fa{اژدها} (ezhdehâ) — a dragon}

\label{lesson:FA-C347-ezhdeha}



\begin{quote}
Before the new one: say the Persian for a myth, then the Persian for a hero.
\end{quote}

\subsection*{You'll want to know: \fa{اژدها}}
\textbf{\fa{اژدها}} (\emph{ezhdehâ}) — \textquotedblleft{}a dragon\textquotedblright{}.

\begin{quote}
\fa{در} \fa{قصه} \fa{یک} \fa{اژدها} \fa{کنار} \fa{رودخانه} \fa{زندگی} \fa{میکرد}.
\end{quote}

\begin{quote}
\emph{Dar qesse yek ezhdehâ kenâr-e rudkhâne zendegi mikard.}
\end{quote}

\begin{quote}
In the tale a dragon lived by the river.
\end{quote}

Written with \fa{ژ}; also said \emph{azhdahâ}.

\begin{cousinweb}
\emph{ezhdehâ} is inherited from Middle Persian \emph{azdahāg} (EtymDB), which continues the Avestan dragon name \emph{aži dahāka}.
\end{cousinweb}

\begin{grammarlens}[title={What you've built}]
One more word for telling, explaining and describing at B1.
\end{grammarlens}

\subsection*{Your turn}
\begin{itemize}
\raggedright
  \item \emph{Say it:} \emph{ezhdehâ}
  \item \emph{Say it:} \emph{ezhdehâ}, once more
  \item \emph{Say it:} a sentence of your own with \emph{ezhdehâ}
  \item \emph{Recall:} say the Persian for a myth, then the Persian for a hero, then say \emph{ezhdehâ} again
\end{itemize}

\begin{tcolorbox}[breakable,colback=teal!4,colframe=teal!35!black,title={Before you move on}]
What does \fa{اژدها} mean? (\textquotedblleft{}A dragon\textquotedblright{}.) Say it once more.
\end{tcolorbox}
\section[\texorpdfstring{\fa{دیو} (div)}{دیو (div)}]{\fa{دیو} (div) — a demon, an ogre}

\label{lesson:FA-C347-div}



\begin{quote}
Before the new one: say the Persian for a hero, then the Persian for a dragon.
\end{quote}

\subsection*{You'll want to know: \fa{دیو}}
\textbf{\fa{دیو}} (\emph{div}) — \textquotedblleft{}a demon, an ogre\textquotedblright{}.

\begin{quote}
\fa{در} \fa{قصه} \fa{دیو} \fa{دختر} \fa{را} \fa{دزدید}.
\end{quote}

\begin{quote}
\emph{Dar qesse div dokhtar râ dozdid.}
\end{quote}

\begin{quote}
In the tale the demon carried off the girl.
\end{quote}

In Persian epics the \emph{div} is a huge, evil being. \emph{Divâne} \textquotedblleft{}mad\textquotedblright{} is built on it.

\begin{cousinweb}
\emph{div} is inherited: EtymDB traces it through Middle Persian \emph{dēw} to Old Persian \emph{daiva-}, \textquotedblleft{}evil spirit\textquotedblright{}, and on to an Indo-European word for a god. Persian turned the old gods into demons.
\end{cousinweb}

\begin{grammarlens}[title={What you've built}]
One more word for telling, explaining and describing at B1.
\end{grammarlens}

\subsection*{Your turn}
\begin{itemize}
\raggedright
  \item \emph{Say it:} \emph{div}
  \item \emph{Say it:} \emph{div}, once more
  \item \emph{Say it:} a sentence of your own with \emph{div}
  \item \emph{Recall:} say the Persian for a hero, then the Persian for a dragon, then say \emph{div} again
\end{itemize}

\begin{tcolorbox}[breakable,colback=teal!4,colframe=teal!35!black,title={Before you move on}]
What does \fa{دیو} mean? (\textquotedblleft{}A demon, an ogre\textquotedblright{}.) Say it once more.
\end{tcolorbox}
\section[(dialogue)]{Words from chapters 340–347 — a review}

\label{lesson:FA-R347-review-340-347}



\begin{quote}
Say the Persian for New Year money and a round of visits.
\end{quote}

\subsection*{Your turn}
Say these aloud:

\begin{itemize}
\raggedright
  \item New Year money, a round of visits, a rite, a custom, wild rue
  \item to throw a party, to serve, a present, a balloon, a candle
  \item to decorate, festive lights, fireworks, a parade, joy
  \item to win, to lose, a spectator, to cheer, a medal
  \item a vow, votive food, halva, saffron rice pudding, rosewater
  \item a tale, once upon a time, there was one and there was none, a storyteller, a narrator
  \item what happened, an account, to narrate, to relate, to remind
  \item a legend, a myth, a hero, a dragon, a demon
\end{itemize}

\begin{tcolorbox}[breakable,colback=teal!4,colframe=teal!35!black,title={Before you move on}]
New Year money? (\textbf{\fa{عیدی}}.) A round of visits? (\textbf{\fa{دید} \fa{و} \fa{بازدید}}.) A rite? (\textbf{\fa{آیین}}.) A custom? (\textbf{\fa{رسم}}.) Wild rue? (\textbf{\fa{اسپند}}.) To throw a party? (\textbf{\fa{مهمانی} \fa{دادن}}.) To serve? (\textbf{\fa{پذیرایی} \fa{کردن}}.) A present? (\textbf{\fa{کادو}}.) A balloon? (\textbf{\fa{بادکنک}}.) A candle? (\textbf{\fa{شمع}}.) To decorate? (\textbf{\fa{تزیین} \fa{کردن}}.) Festive lights? (\textbf{\fa{چراغانی}}.) Fireworks? (\textbf{\fa{آتش} \fa{بازی}}.) A parade? (\textbf{\fa{رژه}}.) Joy? (\textbf{\fa{شادی}}.) To win? (\textbf{\fa{برنده} \fa{شدن}}.) To lose? (\textbf{\fa{باختن}}.) A spectator? (\textbf{\fa{تماشاگر}}.) To cheer? (\textbf{\fa{تشویق} \fa{کردن}}.) A medal? (\textbf{\fa{مدال}}.) A vow? (\textbf{\fa{نذر}}.) Votive food? (\textbf{\fa{نذری}}.) Halva? (\textbf{\fa{حلوا}}.) Saffron rice pudding? (\textbf{\fa{شله} \fa{زرد}}.) Rosewater? (\textbf{\fa{گلاب}}.) A tale? (\textbf{\fa{قصه}}.) Once upon a time? (\textbf{\fa{روزی} \fa{روزگاری}}.) There was one and there was none? (\textbf{\fa{یکی} \fa{بود} \fa{یکی} \fa{نبود}}.) A storyteller? (\textbf{\fa{قصه} \fa{گو}}.) A narrator? (\textbf{\fa{راوی}}.) What happened? (\textbf{\fa{ماجرا}}.) An account? (\textbf{\fa{روایت}}.) To narrate? (\textbf{\fa{روایت} \fa{کردن}}.) To relate? (\textbf{\fa{نقل} \fa{کردن}}.) To remind? (\textbf{\fa{یادآوری} \fa{کردن}}.) A legend? (\textbf{\fa{افسانه}}.) A myth? (\textbf{\fa{اسطوره}}.) A hero? (\textbf{\fa{قهرمان}}.) A dragon? (\textbf{\fa{اژدها}}.) A demon? (\textbf{\fa{دیو}}.)
\end{tcolorbox}
